% SPDX-License-Identifier: Apache-2.0
\section{The Top Is a Program}
\label{sec:e-main}

Section~\ref{sec:e-hierarchy} put units inside units. The outermost one
has to be reached from somewhere, and in an embedded language that
somewhere already exists.

\begin{lstlisting}[caption={A design is a program. The config names the top unit as well as filling it, so \code{main} needs nothing else.},label={lst:main}]
pub trait Config {
    type Top: Module<(), ()>;
    const NAME: &'static str;
    const CLK_HZ: u64;
    fn top() -> Self::Top;
}

fn main() {
    let which = std::env::args().nth(1)
        .unwrap_or_else(|| "fpga".to_string());
    match which.as_str() {
        "asic" => elaborate::<Asic>(),
        _      => elaborate::<Fpga>(),
    }
}
\end{lstlisting}

There is no separate elaboration entry point and no build description
language. Probe~13 is a binary rather than a library for that reason, and
running it prints what it elaborated:

\begin{lstlisting}[style=ascii,caption={Running the design. Each arm of the match is monomorphised separately, so both builds are compiled and either can be asked for.},label={lst:run}]
$ bazel run //experiments/rust_embedding:probe_main -- asic
elaborated config 'asic' at 400 MHz
\end{lstlisting}

\subsection{A build is one item}

Written out, a build was a unit struct, an \code{impl} of the design's
own configuration trait, and an \code{impl Config} whose \code{fn top}
constructed the entire hierarchy field by field. Most of that is
mechanical, and the constructor is the part that grows with the design.

Two changes to the runtime remove it, and a macro removes the rest.
\code{Config::top} builds the top from its \code{Default}, because a
register's default is its reset value and a socket's is its
implementation, so a unit that derives \code{Default} needs no
constructor at all. A configuration is a unit struct that derives
\code{Default} too, because a unit generic over it derives
\code{Default} and the derive asks that of every parameter. Then
\code{config!} writes the struct and both impls:

\begin{lstlisting}[caption={A build is one item. The body is spliced into the design's own configuration impl unchanged, so the macro never parses it.},label={lst:config-macro}]
config! { Fpga: TopConfig for Top<Fpga> {
    type Filter = SixteenTaps;
    const CLK_HZ: u64 = 100_000_000;
} }

config! { Asic: TopConfig for Top<Asic> {
    type Filter = SixtyFourTaps;
    const CLK_HZ: u64 = 400_000_000;
} }
\end{lstlisting}

The example in the appendix went from 81 lines to 55 under this change,
and every line removed was one that repeated a field name.

Three things follow, and the third is the one that matters.

The config gains \code{type Top} and \code{fn top()}, so it says which
design to elaborate as well as what to fill it with. A config is then the
whole of what a build is, which is what the configuration facet was for.

Selecting a build is an ordinary argument. Each arm names a different type
and is compiled separately, so asking for a build that does not exist is a
compile error rather than a missing file.

And the toolchain becomes a library rather than a program that reads
source. \code{elaborate} is a function the design calls, so a project may
wrap it, run two configs and compare them, generate a report, or drive a
simulation from a test, using whatever Rust already offers. None of that
needed designing.
